\documentclass[11pt, a4paper]{article}

\usepackage[T1]{fontenc}
\usepackage[utf8]{inputenc}
\usepackage[spanish]{babel}
\usepackage[a4paper, margin=2.5cm]{geometry}
\usepackage{amsmath, amssymb}
\usepackage{booktabs}
\usepackage{enumitem}
\usepackage{listings}
\usepackage{xcolor}
\usepackage{hyperref}

\definecolor{codigobg}{RGB}{245,245,245}
\definecolor{keyword}{RGB}{0,100,150}
\lstset{language=R, basicstyle=\ttfamily\small, backgroundcolor=\color{codigobg},
  keywordstyle=\color{keyword}\bfseries, frame=single, breaklines=true, showstringspaces=false}

\setlength{\parindent}{0pt}
\pagestyle{plain}

\begin{document}

\begin{center}
{\large\bfseries PONTIFICIA UNIVERSIDAD JAVERIANA}\\[3pt]
Facultad de Ciencias Econ\'omicas y Administrativas\\[2pt]
Estad\'istica 2026-II --- Prof.\ Jaime Polanco Jim\'enez\\[6pt]
{\LARGE\bfseries Problem Set 7}\\[3pt]
{\large Bondad de Ajuste y Regresi\'on Lineal Simple}\\[6pt]
Asignado: Octubre 16 (Sesi\'on 9) \quad$|$\quad Entrega: Octubre 30 (Sesi\'on 10)\\[3pt]
\textit{Plazo extendido: 2 semanas}
\end{center}

\vspace{6pt}
\rule{\linewidth}{0.5pt}
\vspace{2pt}
{\centering\small Repositorio: \url{https://github.com/polanco-jaime/Curso-Estadistica-2026-3}\par}
\vspace{6pt}

\textbf{Instrucciones.} Este Problem Set es individual y se resuelve en R sobre
microdatos ICFES. La entrega es digital al inicio de la clase indicada.
El c\'odigo R debe incluirse.

\vspace{10pt}

\section*{Enunciado}

En este Problem Set evaluar\'a si la distribuci\'on de niveles de ingl\'es
seg\'un el MCER es uniforme entre estudiantes de Negocios Internacionales (NI),
y estimar\'a un modelo de regresi\'on lineal simple para predecir el rendimiento
en Saber Pro a partir de Saber 11.

\begin{enumerate}[label=\arabic*)]
  \item \textbf{Prueba $\chi^2$ de bondad de ajuste:} ¿los niveles de ingl\'es
        seg\'un el Marco Com\'un Europeo de Referencia (MCER) se distribuyen
        uniformemente entre los estudiantes de NI?
        \begin{enumerate}[label=\alph*)]
          \item Construya una tabla de frecuencias observadas para los niveles
                MCER (A$-$, A1, A2, B1, B2) en la variable
                \texttt{MOD\_INGLES\_DESEM\_MCER}.
          \item Plantee las hip\'otesis:
                \[
                  H_0: p_1 = p_2 = \cdots = p_5 = \frac{1}{5}
                  \quad \text{vs.} \quad
                  H_1: \text{al menos una } p_i \text{ difiere de } \frac{1}{5}
                \]
          \item Calcule las frecuencias esperadas bajo $H_0$ (distribuci\'on
                uniforme).
          \item Aplique la prueba $\chi^2$ de bondad de ajuste usando
                \texttt{chisq.test()} y reporte el estad\'istico $\chi^2$,
                grados de libertad y $p$-valor.
          \item ¿Rechaza $H_0$ al nivel $\alpha = 0.05$? Interprete el
                resultado: ¿qu\'e niveles MCER est\'an sobrerrepresentados
                o subrepresentados?
        \end{enumerate}

  \item \textbf{Regresi\'on lineal simple (OLS):} estime un modelo para predecir
        el puntaje de ingl\'es en Saber Pro (\texttt{MOD\_INGLES\_PUNT}) a partir
        del puntaje de ingl\'es en Saber 11 (\texttt{PUNT\_INGLES}).
        \begin{enumerate}[label=\alph*)]
          \item Prepare los datos: cruce las bases de Saber 11 y Saber Pro
                usando el dataset \texttt{ph-jabri.ICFES.Saber11\_SaberPro}.
                Filtre estudiantes de NI con datos completos en ambas pruebas.
          \item Grafique un diagrama de dispersi\'on (scatterplot) de
                \texttt{PUNT\_INGLES} (eje $x$) vs.\
                \texttt{MOD\_INGLES\_PUNT} (eje $y$). ¿Observa una relaci\'on
                lineal?
          \item Estime el modelo OLS:
                \[
                  \texttt{MOD\_INGLES\_PUNT}_i = \beta_0 + \beta_1
                  \texttt{PUNT\_INGLES}_i + \varepsilon_i
                \]
                usando \texttt{lm()} y presente el \texttt{summary()} completo.
          \item Interprete los coeficientes $\hat{\beta}_0$ y $\hat{\beta}_1$:
                \begin{itemize}
                  \item ¿Qu\'e representa $\hat{\beta}_0$ (intercepto)?
                  \item ¿Qu\'e representa $\hat{\beta}_1$ (pendiente)? Si el
                        puntaje de Saber 11 aumenta en 10 puntos, ¿cu\'anto
                        se espera que aumente el puntaje de Saber Pro?
                \end{itemize}
        \end{enumerate}

  \item \textbf{Diagn\'ostico de residuos:} eval\'ue la calidad del ajuste del
        modelo OLS.
        \begin{enumerate}[label=\alph*)]
          \item Calcule los residuos $\hat{\varepsilon}_i = y_i - \hat{y}_i$
                y grafique:
                \begin{itemize}
                  \item Residuos vs.\ valores ajustados $\hat{y}_i$
                  \item Q-Q plot de residuos
                \end{itemize}
          \item Interprete los gr\'aficos: ¿los residuos parecen cumplir los
                supuestos de linealidad, homocedasticidad y normalidad?
          \item Identifique las 3 observaciones con mayor residuo absoluto
                (outliers potenciales). ¿Qu\'e caracter\'isticas tienen?
        \end{enumerate}

  \item \textbf{Bondad de ajuste ($R^2$):} eval\'ue el poder explicativo del
        modelo.
        \begin{enumerate}[label=\alph*)]
          \item Reporte el coeficiente de determinaci\'on $R^2$ y el $R^2$
                ajustado.
          \item Interprete $R^2$: ¿qu\'e porcentaje de la variabilidad en el
                puntaje de Saber Pro es explicado por el puntaje de Saber 11?
          \item ¿El modelo tiene un buen ajuste? Justifique.
        \end{enumerate}

  \item \textbf{Predicci\'on:} use el modelo estimado para predecir el puntaje
        de Saber Pro de un estudiante de NI que obtuvo 75 puntos en ingl\'es
        en Saber 11.
        \begin{enumerate}[label=\alph*)]
          \item Calcule la predicci\'on puntual $\hat{y}$.
          \item Calcule un intervalo de confianza al 95\% para la predicci\'on
                usando \texttt{predict()} con \texttt{interval = "prediction"}.
          \item Interprete el intervalo: ¿en qu\'e rango se espera que est\'e
                el puntaje de Saber Pro con 95\% de confianza?
        \end{enumerate}

  \item \textbf{Conclusiones:} sintetice los hallazgos:
        \begin{enumerate}[label=\alph*)]
          \item ¿La distribuci\'on de niveles MCER es uniforme en NI?
          \item ¿El puntaje de Saber 11 es un buen predictor del puntaje de
                Saber Pro en ingl\'es?
          \item ¿Qu\'e limitaciones tiene el modelo OLS simple estimado?
        \end{enumerate}
\end{enumerate}

\vspace{10pt}
\rule{\linewidth}{0.4pt}
\begin{center}
{\small Cada entrega completa suma $+0.0625$ a la nota final.}
\end{center}

\end{document}
